\documentclass[12pt, letterpaper]{article}

%-------Packages---------
\usepackage{newpxtext,newpxmath} % Palatino-like text & math (modern)
\usepackage{bboldx}
\usepackage{mathtools}
\usepackage{tikz}
\usetikzlibrary{calc,intersections}
\usepackage{tikz-cd}
\usepackage[all,arc]{xy}
\usepackage{graphicx}

\usepackage{enumerate}
\usepackage[dvipsnames]{xcolor}
\usepackage{float}
\usepackage[left=1in,top=1in,right=1in,bottom=1in]{geometry}
\usepackage[nocheck]{fancyhdr}
\usepackage{multicol}
\usepackage{setspace}

\usepackage{dsfont}
\usepackage{mathrsfs}

\usepackage{tcolorbox}
\tcbuselibrary{theorems,skins,breakable}

\usepackage{xparse}
\usepackage{import}
\usepackage[utf8]{inputenc}

\usepackage{hyperref}
\usepackage{cleveref}

\usepackage{xcolor, ifthen, tcolorbox}
\tcbuselibrary{theorems,skins,breakable}
% Theorem System
% The following environments are provided:
%   New Problem:        \begin{problem} ...
%   New Question Header:   \qh (then followed by subproblems)
%   Subproblem:     \begin{subproblem} ...
%   Problem Proof:  \begin{problemproof} ...
%   Proof:          \\begin{pf}{color} ...
%   Remark:         \rmk (sentence), \rmkb (block)

% Define custom colors
\definecolor{thmscol}{HTML}{F4565B} %For Problems
\definecolor{lemscol}{HTML}{FE844C} %For Lemmes
\definecolor{prpscol}{HTML}{00A7EF} %For proposition
\definecolor{corscol}{HTML}{23DAFF} %For corrolaries
\definecolor{rmkscol}{HTML}{6DEEC5} %For remarks
\definecolor{defscol}{HTML}{18C991} %For definitions
\definecolor{exmscol}{HTML}{7DB800} %For examples
\definecolor{clmscol}{HTML}{165c58} %For claims

%Change between dark(black)/light(white) mode
\newcommand{\colormode}{white}
\ifthenelse{\equal{\colormode}{white}}
      {\newcommand{\TextColor}{black}}
      {\newcommand{\TextColor}{white}}
\newcommand{\pbcolormain}{thmscol!80!\colormode}
\newcommand{\pbcolorsecond}{thmscol!38!\colormode}


% ============================
% Problem Counter
% ============================
\newcounter{subproblemcounter}
\newcounter{problemcounter}

\newcommand{\q}{
    \setcounter{subproblemcounter}{0}%
    \stepcounter{problemcounter} % Increment problem counter
    \color{\TextColor}Problem \arabic{problemcounter}: % Display the problem counter
}

\newcommand{\stepsub}{
    \stepcounter{subproblemcounter}%
    \color{\TextColor}(\alph{subproblemcounter})
}
% ============================
% Problem
% ============================
\newtcolorbox{pb}[1]{
  enhanced,
  frame hidden,
  titlerule=0mm,
  toptitle=1mm,
  bottomtitle=1mm,
  fonttitle=\bfseries\large,
  coltitle=black,
  colbacktitle=\pbcolormain,
  colback=\pbcolorsecond,
  title={\q #1},
  coltext=\TextColor
}

\newtcolorbox{spb}[1]{
  enhanced,
  frame hidden,
  titlerule=0mm,
  toptitle=1mm,
  bottomtitle=1mm,
  fonttitle=\bfseries\large,
  coltitle=black,
  colbacktitle=\pbcolormain,
  colback=\pbcolorsecond,
  title={\stepsub #1},
  coltext=\TextColor,
  attach boxed title to top left={xshift=3mm,yshift=-2mm},
  boxed title style={
    colback=\pbcolormain,
    colframe=\pbcolormain,
  }
}

\newtcolorbox{pbh}[1]{
    enhanced,
    frame hidden,
    titlerule=0mm,
    toptitle=1mm,
    bottomtitle=1mm,
    fonttitle=\bfseries\large,
    coltitle=black,
    colbacktitle=\pbcolormain,
    colback=\pbcolorsecond,
    title={\q #1},
    boxrule=0pt,
    interior hidden,
    attach boxed title to top left={xshift=3mm,yshift=-2mm},
    boxed title style={
        colback=\pbcolormain,
        colframe=\pbcolormain,
    }
}

\newenvironment{problem}[1]{%
  \newpage
  \begin{pb}{#1}
    }{
  \end{pb}
}

\newenvironment{subproblem}[1]{%
  \begin{spb}{#1}
    }{
  \end{spb}
}

\newenvironment{problemproof}{%
    \begin{pf}{\pbcolormain}
        }{
    \end{pf}
}

\newcommand{\qh}[1][]{
    \newpage
    \begin{pbh}{#1}\end{pbh} % Display the problem counter
}

% ============================
% Claim
% ============================

\newtcbtheorem[number within=problemcounter]{claim}{Claim}
{
    enhanced,
    frame hidden,
    titlerule=0mm,
    toptitle=1mm,
    bottomtitle=1mm,
    fonttitle=\bfseries\large,
    coltitle=black,
    colbacktitle=clmscol!50!\colormode,
    colback=clmscol!20!\colormode,
}{clm}

\NewDocumentCommand{\clm}{m+m}{
    \begin{claim*}{#1}{}
        #2
    \end{claim*}
}

\newenvironment{clmpf}{
	{\noindent{\it \textbf{Proof.}}
	\setlength{\parindent}{0pt}
	}
	\tcolorbox[blanker,breakable,left=5mm,parbox=false,
    before upper={\parindent15pt},
    after skip=10pt,
	borderline west={1mm}{0pt}{clmscol!50!\colormode}]
}{
    \textcolor{clmscol!50!\colormode}{\hbox{}\nobreak\hfill$\blacksquare$}
    \endtcolorbox
}

\NewDocumentCommand{\clmp}{m+m+m}{
    \begin{claim*}{#1}{}
        #2
    \end{claim*}

    \begin{clmpf}
        #3
    \end{clmpf}
}


% ============================
% Proof
% ============================

\newenvironment{pf}[1]{%
  \def\pfcolor{#1}% Store the color in a macro
  \noindent{\it \textbf{Proof.}}%
  \tcolorbox[blanker,breakable,left=5mm,parbox=false,%
    before upper={\parindent15pt},%
    after skip=10pt,%
    borderline west={1mm}{0pt}{\pfcolor},%
    coltext=\TextColor
  ]%
  \setlength{\parindent}{0pt}%
}{%
  \textcolor{\pfcolor}{\hbox{}\nobreak\hfill$\blacksquare$}%
  \endtcolorbox%
}

\newtcolorbox{solution}{
  blanker,
  breakable,
  left=5mm,
  parbox=false,%
  before upper={\parindent15pt},%
  after skip=10pt,%
  borderline west={1mm}{0pt}{\pbcolormain},%
  coltext=\TextColor
}


% ============================
% Remark
% ============================


\NewDocumentCommand{\rmk}{+m}{
    {\it \color{rmkscol!100!white}#1}
}

\newenvironment{remark}{
    \par
    \vspace{5pt}
    \begin{minipage}{\textwidth}
        {\par\noindent{\textbf{Remark.}}}
        \tcolorbox[blanker,breakable,left=5mm,
        before skip=10pt,after skip=10pt,
        borderline west={1mm}{0pt}{rmkscol!80!\colormode},
        coltext=\TextColor
        ]
}{
        \endtcolorbox
    \end{minipage}
    \vspace{5pt}
}

\NewDocumentCommand{\rmkb}{+m}{
    \begin{remark}
        #1
    \end{remark}
}


% Define a new command for problems
\renewcommand{\headrule}{\color{\TextColor}\hrulefill}
\newcommand{\C}{\mathbb{C}}
\newcommand{\R}{\mathbb{R}}
\newcommand{\Q}{\mathbb{Q}}
\newcommand{\Z}{\mathbb{Z}}
\newcommand{\N}{\mathbb{N}}
\newcommand{\p}{\mathfrak{p}}
\newcommand{\F}{\mathbb{F}}
\newcommand{\contr}{\Rightarrow \Leftarrow}
\newcommand{\s}{\(\,\)}
\renewcommand{\t}[1]{\text{#1}}
\newcommand{\Spec}{\mathrm{Spec}}
\newcommand{\Ass}{\mathrm{Ass}}
\newcommand{\Supp}{\mathrm{Supp}}
\newcommand{\Ann}{\mathrm{Ann}}
\newcommand{\mf}[1]{\mathfrak{#1}}
\newcommand{\codim}{\mathrm{codim}}

% Meta data
\setstretch{1.2}

\begin{document}
\color{\TextColor}
\pagecolor{\colormode}
\setlength{\parindent}{0pt}
\pagestyle{fancy}
\fancyhf{}
\fancyhead[LOE]{\color{\TextColor}\slshape{\textbf{Vincent Ferrara}}}
\fancyhead[ROE]{\color{\TextColor}HW10 --- \thepage}
\fancyhead[C]{\color{\TextColor}Math 614}

\qh
\addtocounter{subproblemcounter}{1}
\begin{subproblem}{}
    Theorem. Let $R$ be a domain which is a finite type $k$-algebra for a field $k$. Then for
any ideal $I \subset R$, we have
\[\dim R = \dim R \diagup I + \codim I\]
\end{subproblem}
\begin{problemproof}
    Let $\p \in \Spec(R)$ give a maximal chain with $\p$
    \[(0)=\p_0 \subsetneq \p_1 \subsetneq \dots \subsetneq \p_r=\p \subsetneq \dots \subsetneq \p_n\]
    and by (a) we get that $n=\dim R$ for all of $\p$'s maximal chains.

    Consider $\p_r \subsetneq \dots \subsetneq \p_n$

    This can also be considered a chain in $R \diagup \p$.

    Thus, $\dim R \diagup \p \geq \dim R - \codim \p$.

    Now let $0 \subsetneq \mf{q}_1 \subsetneq \dots \subsetneq \mf{q}_r$ be a max chain in $R \diagup \p$ s.t. $\dim R \diagup \p = r$

    This corresponds to a chain $\p \subsetneq \p_1 \subsetneq \dots \subsetneq \p_r$ in $R$.

    Now extend this chain downwards s.t. we achieve $\p$'s codimension. Thus, we get a chain of length $r+\codim \p$.

    Thus, $\codim \p + r \leq \dim R$.

    Therefore, we get that $\dim R = \dim R \diagup \p + \codim \p$ for $\p$ prime.

    Now let $I$ be an ideal in $R$. Since $R$ is noetherian we know that $I$ has finitly many minimal prime ideals over it.

    Thus, $\dim R \diagup I = \max \{\dim R \diagup \p \mid \p \t{ is minimal over } I\}$

    Also, by definition of codimension we have $\codim I = \{\codim \p \mid \p \t{ is minimal over } I\}$.

    Thus, applying the lemma above to each $\p$ we get that $\dim R \diagup \p = \dim R - \codim \p$.

    Thus, $\dim R \diagup I = \max \{\dim R \diagup \p \mid \p \t{ is minimal over } I\}$

    $=\max \{\dim R - \codim  \p \mid \p \t{ is minimal over } I\} = \dim R - \min \{\codim \p \mid \p \t{ is minimal over } I\}$.

    Thus, $\dim R \diagup I = \dim R - \codim I$.
\end{problemproof}

\qh
\begin{subproblem}{}
    Compute the dimension of $\Q[x,y,\sqrt{x^7 + y^2}]$.
\end{subproblem}
\begin{problemproof}
    Since $R=\Q[x,y,\sqrt{x^7+y^2}] \cong \Q[x,y,z] \diagup (z^2-x^7-y^2)$. This implies 
    $R$ is a finite type $\Q$-algebra. Thus, by problem 1 we get that $\dim \Q[x,y,z] = \dim R + \codim (z^2-x^7-y^2)$.

    Since $(z^2-x^7-y^2)$ is irreducible this implies it is prime, and this also implies that no ideal but the zero ideal is contained in this. Therefore, since $\Q[x,y,z]$
    is a domain this implies that $(z^2-x^7-y^2)$ has codimension $1$.

    Therefore, $\dim R = 3-1=2$.
\end{problemproof}

\begin{subproblem}{}
    Compute the dimension of the ring $R=\Z[x, y, z, w]/(x^3 + yzw + w^5)$.
\end{subproblem}
\begin{problemproof}
    Since $(x^3+yzw + w^5)$ is irreducible this implies it is prime. Thus, by the PIT we get that $\codim (x^3+yzw+w^5) \leq 1$, but since we are in a domain its codimension is one. Thus, shown in class we have that $\dim \Z[x,y,z,w] - \codim I \geq \dim R \implies 4 \geq \dim R$.

    Also since $(x^3+yzw+w^5) \subsetneq (x,w) \subsetneq (2,x,w) \subsetneq (2,x,y,w) \subsetneq (2,x,y,z,w)$ is a chain. This implies that $4 \leq \dim R$. Therefore, this implies $\dim R = 4$.
\end{problemproof}

\begin{subproblem}{}
    Let $f \in \Q[x,y,z]$ be a nonzero polynomial. Compute the dimension of the ring $\Q[x,y,z]_f$.
\end{subproblem}
\begin{problemproof}
    Since $\Q[x,y,z]_f \cong \Q[x,y,z,t] \diagup (tf-1)$.

    Again by problem 1 since $(tf-1)$ is irreducible thus prime by PIT its codimension is at most one, but since we are in a domain this implies it has codimension one. Thus, $\dim \Q[x,y,z]_f = \dim \Q[x,y,z,t] - \codim (tf-1) = 3$.
\end{problemproof}

\begin{subproblem}{}
    Let $R = \C[x,y,z] \diagup (xy,xz)$. Compute the dimension of $R_{\bar{x}}$.
\end{subproblem}
\begin{problemproof}
    Since $\R_{\bar{x}} \cong \C[x,y,z,t] \diagup (xy,xz,tx-1)$. 
    
    Also since $x$ is now a unit in $\C[x,y,z,t] \diagup (xy,xz,tx-1)$ this implies we can cancel the $x$'s thus, $\C[x,y,z,t] \diagup (xy,xz,tx-1) = \C[x,y,z,t] \diagup (y,z,tx-1) \cong \C[x,t] \diagup (tx-1)$. Then from a similar proof to $(c)$ we get that $\dim \C[x,t] \diagup (tx-1) = 1$ since $\dim \C[x,t]=2$ and $\codim (tx-1)=1$.
\end{problemproof}

\begin{subproblem}{}
    Compute the codimension of the ideal $(x_1, \dots, x_m) \subset k[x_1,\dots , x_n]$, where $m \leq n$ and
$k$ is a field.
\end{subproblem}
\begin{problemproof}
    From problem one since $\k[x_1,\dots,x_n]$ is a finite type $k$ algebra this implies $\codim (x_1,\dots,x_m) = m$ since $\dim \k[x_1,\dots,x_n] = n$ and $\dim \k[x_1,\dots,x_n] \diagup (x_1,\dots,x_m) = n -m$.
\end{problemproof}

\begin{problem}{}
    This problem is about UFDs from the point of Noetherian rings and dimension theory.
\end{problem}
\addtocounter{subproblemcounter}{1}
\begin{subproblem}{}
    Prove that a domain $R$ is a UFD if and only if the following conditions holds:
    \begin{enumerate}[i.]
        \item For every nonzero, noninvertible element $x \in R$ there exists a factorization
$x = \pi_1 \dots \pi_r$ where $\pi_1, \dots , \pi_r \in R$ are irreducible.
        \item Every irreducible element in $R$ is prime.
    \end{enumerate}
\end{subproblem}
\begin{problemproof}
    Assume $R$ is a UFD.

    $i.$ is immediate from the definition of UFD.

    Let $x \in R$ be irreducible.

    Let $x \mid ab$. Then $ab=xr$ s.t. $r \in R$.

    Then factor $a,b,c$ into irreducibles. By the uniqueness of factorizations this implies that $x$ must appear
    among the irreducible factors of either $a$ or $b$. Thus, $x \mid a$ or $x \mid b$.

    Therefore, $x$ is prime.

    Assume $i.$ and $ii.$

    First existence of factorization exits from $i$.

    Suppose $x\in R$ has two factorizations $x=p_1\dots p_n = q_1 \dots q_m$.

    Since $p_1 \mid q_1 \dots q_m$ by $ii.$ this implies $p_1 \mid q_j$. Thus, $q_j = u p_1$ for some unit $u$.

    Therefore, you can do this over and over which implies $n=m$ and $p_i = u_j q_j$ for units $u_j$.

    Therefore, factorizations are unique up to units.

    Thus, $R$ is a UFD.
\end{problemproof}

\begin{subproblem}{}
    Prove that if $R$ is a noetherian domain, then for every nonzero, noninvertible element
$x \in R$, there exists a factorization $x = \pi_1 \dots \pi_r$ where $\pi_1, \dots , \pi_r \in R$ are irreducible
elements. In particular, combined with (b), we see that a noetherian domain is a UFD
if and only if every irreducible element is prime.
\end{subproblem}
\begin{problemproof}
    Assume $R$ is a noetherian domain.

    Let $S = \{x \in R \setminus \{0\} \mid x \t{ is not a unit, and $x$ does not factor into irreducibles}\}$.

    Assume for contradiction that $S \neq \emptyset$.

    Then consider $\{(x) \mid x \in S\} \neq \emptyset$. Since $R$ is noetherian this implies that there exists a maximal element $(y)$ in this set.

    Since $y \in S$ this implies $y$ is not irreducible thus $y=ab$ s.t. $a,b \in R$ and are not units.

    Since $(x) \subsetneq (a),(b)$. This implies $a,b$ do factor into irreducibles since $(x)$ is maximal thus, $a,b \notin S$.

    Thus, $a=p_1\dots p_n$ and $b=q_1 \dots q_m$, but this implies $x=p_1\dots p_n q_1 \dots q_m$ which is a contradiction. Therefore, $S = \emptyset$.

    Therefore, all noninvertible elements factor into irreducibles.
\end{problemproof}

\begin{subproblem}{}
    Prove that if $R$ is a noetherian domain, then $R$ is a UFD if and only if every prime
ideal $\p \in \Spec(R)$ with $\codim \p = 1$ is a principal ideal.
\end{subproblem}
\begin{problemproof}
    Assume $R$ is a noetherian domain and a UFD.

    Let $\p \in \Spec(R)$ s.t. $\codim \p = 1$. Let $\p=(a,c_1,\dots,c_n)$

    Then consider $(a) \subsetneq \p$. Since $R$ is a domain this implies that $(a)$ is not prime since if it was then we would have $(0) \subsetneq (a) \subsetneq \p$. Which would implies $\p$ has codimension 2.

    Thus, $(a)$ is not prime. Thus, $a$ is not irreducible this implies $a=p_1 \dots p_n$ s.t. $p_i$ are irreducible. Consider $(p_1)$ which is prime since we are in a UFD. This implies that $(0) \subsetneq (p_1)$ is a chain. Thus,
    $(p_1)$ has codimension 1 by the PIT, but since $(p_1) \subset \p$ this implies that $(p_1) = \p$. Therefore, $\p$ is a principal ideal.
\end{problemproof}

\qh
\begin{subproblem}{}
    
\end{subproblem}
\begin{problemproof}
    Let $f \in k\{\{t\}\}$ s.t. $f=\sum_{i=i_0}^\infty a_i t^{i/n}$ for some $i_0 \in \Z$ and $n \in \N$, and $a_i \in k$ and $a_{i_0} \neq 0$.

    Suppose $f$ also is equal to $\sum_{j=j_0}^\infty b_j t^{j/m}$ for some $j_0 \in \Z$ and $m \in \N$, and $b_j \in k$ and $b_{j_0} \neq 0$.

    Let $N = \mathrm{lcm}(n,m)$. Then this implies both can be written as s.t. $i/n = (iN/n)/N$ and same for $j/m$.
    \[f=\sum_{r=r_0}^\infty c_r t^{r/N} = \sum_{r=r_0}^\infty d_r t^{r/N}\]

    This implies $c_r = d_r$ for all $r$.

    Therefore, this implies that the first non zero term for $c_r$ has power on $t$ eaul to $i/n=(iN/n)/N$ since these are the same series this implies that $(iN/n)/N=(jN/m)/M \implies i/n=j/m$.

    Thus, the valuation map is well-defined.

    Let $f \in k\{\{t\}\}$ s.t. $v(f)=\infty$ this implies that there is no nonzero terms in $f$. Thus, $f=0$.

    If $f=0$ then this implies that there is no minimal nonzero term thus $v(f)=\infty$.

    Let $f,g \in k\{\{t\}\}$ s.t.
    \begin{align}
        f = \sum_{i=i_0}a_i t^{i/n}, \quad a_{i_0} \neq 0\\
        g = \sum_{j=j_0}b_j t^{j/m}, \quad b_{j_0}  \neq 0
    \end{align}
    Then $fg = \sum_{k=mi_0+nj_0} c_k t^{k/(nm)}$ s.t. $c_k = \sum_{mi+nj=k}a_ib_j$.

    Thus, looking at the smallest possible exponent we get $c_{mi_0+nj_0}=a_{i_0}b_{j_0} \neq 0$ since $a_{i_0} \neq 0 \neq b_{j_0}$.

    Therefore, $v(fg)=(mi_0+nj_0)/(nm) = i_0/n + j_0/m = v(f)+v(g)$.

    If one of $f,g$ are zero then $v(fg)=\infty = v(f)+v(g)$.

    Let $f,g \in k\{\{t\}\}$.

    If $g = 0$ then $v(f+g) = v(f) \geq \min \{v(f), \infty\}= v(f)$.

    So assume $f,g \neq 0$.

    First choose some common denominator for the exponents in the series call this $n$.

    Thus, $f=\sum_{i=i_0}a_it^{i/n}$ and $g=\sum_{j=j_0}b_jt^{j/n}$.

    Let $f+g=\sum_{k=l_0} c_k t^{k/n}$

    If $f+g = 0$ then $v(f+g)= \infty \geq \min\{v(f),v(g)\}$

    Assume $f+g \neq 0$ then if $v(g) < v(f)$ then $v(f+g)=j_0/n  \geq v(g)$.

    If $v(g)=v(f)$ Then

    if $c_{i_0} \neq 0$ then $v(f+g)= i_0/n \geq i_0/n$

    if $c_{i_0} = 0$ then this implies that the first nonzero term has a greater exponent then $i_0$. Thus, $v(f+g)=l_0 > i_0=j_0$.

    Therefore, $v(f+g) \geq \min \{v(f),v(g)\}$

    Therefore, $v$ is a valuation map.
\end{problemproof}

\begin{subproblem}{}
    Let $k$ be a field and let $k(x, y)$ be the fraction field of $k[x, y]$. Show that there is a
valuation $v : k(x, y) \to \Z^{\oplus 2} \cup \{\infty\}$, where $\Z^{\oplus 2}$ is given the lexicographical order, which
is uniquely determined by $v(x^ay^b) = (a, b) for (a, b) \in \Z^{\oplus 2}$.
\end{subproblem}
\begin{problemproof}
    Define $w : k(x,y) \to \Z^{\oplus 2} \cup \infty$ s.t. $w(0) = \infty$
    and for $f \in k[x,y]$ s.t. $f=\sum c_{i,j}x^iy^j$ that $w(f)=(a,b)$ s.t. $(a,b)$ is the min element in $\{(i,j) \in \N^2 \mid c_{i,j} \neq 0\}$.

    Now for $f \in k(x,y)$ s.t. $f=g/h$ s.t. $g,h \in k[x,y]$ and $h \neq 0$ that $w(f)=w(g)-w(h)$.

    Multiplicativity on $k[x,y]$:

    Let $f,g \in k[x,y]$ then consider $w(fg)$ this implies the new minimal nonzero term are just the minimal terms of $f,g$ multiplied together. Thus, if $w(f)=(a,b)$ and $w(g)=(i,j)$ then $w(f+g)=(a+i,b+j) = w(f)+w(g)$.
    
    Well-Defined on $k(x,y)$:

    Suppose $f/g = h/e$. This implies $fe=hg \neq 0$. Thus, by above we have that $w(fe)=w(hg) \implies w(f)-w(g)=w(h)-w(e)$. Thus, this is well-defined.
    
    Multiplicativity on $k(x,y)$:

    $w(f/g \cdot h/e) = w(fh) - w(ge) = (w(f)-w(g)) + (w(h)-w(e)) = w(f/g) + w(h/e)$.

    Inequality:

    $w(f/g + h/e) = w((fe+hg)/ge) = w(fe+hg) - w(ge)$.

    Consider $w(fe+hg)$. Assume $fe$ or $hg$ are not zero since if $fe$ is then we get $w(fe+hg)=w(fe) \geq w(fe)$. Same with if $fe+hg=0$ then $w(fe+hg)=\infty \geq \min\{w(fe),w(hg)\}$.
    
    If $w(fe) < w(hg)$ then $w(fe+hg) = w(fe) \geq \min\{w(fe),w(gh)\}$.

    If $w(fe)=w(hg)$ and the minimal term with exponent $w(fe)$ is not zero. Then $w(fe+hg) =w(fe) \geq w(fe)$.

    If the term with exponent $w(fe)$ is zero. This implies the minimal nonzero term is bigger. Thus, $w(fe+hg) \geq \min w(fe)$.

    Thus, $w(f/g+h/e) \geq \min\{w(fe)+w(hg)\} - w(ge) = \min\{w(f)+w(e),w(h)+w(g)\} - w(g) - w(e) \geq \min\{w(f)-w(g), w(g)-w(e)\} = \min\{w(f/g),w(g/e)\}$.

    Also, if $w(a) = \infty \iff a = 0$ by definition of $w$.

    Now for $x^ay^b$ for $(a,b) \in \N^2$ we get that $v(x^ay^b) = (a,b) = w(x^a,y^b)$.

    So for general $a,b \in \Z^{\oplus 2}$ pick $A,B$ large enough s.t. $a+A \geq 0$ and $b+B \geq 0$ thus $x^ay^b = (x^{a+A}y^{b+B})/(x^Ay^B)$.

    Thus, $w(x^ay^b)=w(x^{a+A}y^{b+B})-w(x^Ay^B)=(a,b)=v(x^ay^b)$.

    uniqueness:

    Let $v$ be a valuation on $k(x,y)$ to $\Z^{\oplus 2} \cup \{\infty\}$ with the property $v(x^ay^b)=(a,b)$.

    Let $f \in k[x,y]$ s.t. $f \neq 0$.

    \[f = \sum c_{i,j}x^iy^j\]
    let $(a_0,b_0)$ be the exponents on the minimal degree term with lexicographical order.

    This implies that $f=x^{a_0}y^{b_0}c_{(a_0,b_0)}(1+c_{(a_0,b_0)}^{-1}h(x,y))$ s.t. $h(x,y) \in k[x,y]$.

    Consider $v(1+c_{(a_0,b_0)}^{-1}h(x,y)) \geq v(1)=0$, also $v(h) > (0,0)$ since is either zero on a polynomial with a zero constant term since we factored out all the lower degree terms.

    Thus, $v(h) > (0,0)$.

    Consider $1=(1+h)-h$.

    $(0,0)=v(1)=v(1+h-h) \geq \min\{v(1+h),v(h)\}$ but since $v(h) > (0,0)$ we know that $(0,0) \geq v(1+h)$.

    Also since $v(1+h) \leq \min\{v(h),v(1)\} = (0,0)$. This implies $v(1+h)=v(0,0)$.

    Thus, $v(1+c_{(a_0,b_0)}^{-1}h) =(0,0)$ also since $v(c_{(a_0,b_0)}^{-1}h) = (0,0) + v(h)$.

    So, do the same proof with that therefore, applying $v$ to $f$ we get that.

    $v(f=v(x^{a_0}y^{b_0}c_{(a_0,b_0)}(1+c_{(a_0,b_0)}^{-1}h(x,y))))=v(x^{a_0}y^{b_0})+v(c_{(a_0,b_0)})+v(1+c_{(a_0,b_0)}^{-1}h)=(a,b)+(0,0)+(0,0)=(a+b)$.

    Therefore, $v(f)=w(f)$.

    Thus, for $f \in k(x,y)$ s.t. $f \neq 0$.

    Then consider $v(1)=v(ff^{-1}) \implies 0 = v(f) + v(f^{-1}) \implies v(f)=-v(f^{-1})$.

    Thus, consider $f/g \in k(x,y)$.

    This implies $v(f/g)=v(fg^{-1})=v(f)-v(g)=w(f)-w(g)=w(f/g)$. Therefore, $v=w$.
\end{problemproof}

\begin{problem}{}
Let $v' : K \to \Gamma' \cup \{\infty\}$ be any valuation on $K$ such that $R = \{x \in K \mid v'(x) \geq 0\}$ is the
valuation ring of $v'$. Prove that there is an injective order preserving group homomorphism
$\phi : \Gamma \to\Gamma'$ such that $v'(x) = \phi(v(x))$ for all $x \in K^\times$.
\end{problem}
\begin{problemproof}
    Define $\phi : \Gamma \to \Gamma'$ s.t. $\phi([x])= v'(x)$.

    Well-Defined:

    If $[x]=[y]$ in $\Gamma$ then $x/y, y/x \in R$ thus, $v(x/y),v(y/x) \geq 0$.

    Thus, $v'(x)=v'(yxy^{-1})=v'(y)+v'(xy^{-1}) \geq v'(y)$ and $v'(y)=v'(xyx^{-1})=v'(x)+v'(yx^{-1}) \geq v'(x)$.

    Therefore, $\phi([x])=v'(x)=v'(y)=v([y])$.

    homomorphism:

    Consider $\phi([x]+[y])=\phi([xy]) = v'(xy) = v'(x)+v'(y)=\varphi([x])+\varphi([y])$.

    Order:

    Let $[x] \leq [y]$ this implies $y/x \in R \implies v'(y/x) \geq 0$.

    Thus, $\phi([y])=v'(y)=v'(xyx^{-1}) = v'(x) + v'(yx^{-1}) \geq v'(x) = \phi([x])$.
    
    Injective:

    Let $[x] \in \Gamma$ s.t. $\phi([x]) = 0$.

    This implies $v'(x)=0$ thus $0=-v'(x)=v'(x^{-1})$. This implies that $x,x^{-1} \in R$. Thus, $x/1 \in R$ and $1/x \in R$. This implies that $[1]=[x]$. Thus, $\phi([x])=\phi([1])=v'(1)=0$.

    Thus, injective.

    Also since $v(x)=[x]$ this implies $\phi(v(x))=\phi([x])=v'(x)$.
\end{problemproof}
\end{document}